\textbf{Phần 1. Trắc nghiệm lựa chọn}
\setcounterpageref{ex}{0}

\begin{ex} %[3D4N24]
Muốn gấp một số lên nhiều lần, ta lấy
\choice
{Số đó cộng với số lần}
{Số đó chia cho số lần}
{Số đó trừ đi số lần}
{\True Số đó nhân với số lần}
\loigiai{
Muốn gấp một số lên nhiều lần, ta lấy số đó nhân với số lần.\\
\textbf{Đáp án của câu hỏi là D}
}
\end{ex}

\begin{ex} %[3D4N24]
Gấp 6 lên 5 lần ta được
\choice
{$11$}
{\True $30$}
{$35$}
{$25$}
\loigiai{
$6 \times 5 = 30$.\\
\textbf{Đáp án của câu hỏi là B}
}
\end{ex}

\begin{ex} %[3D4H24]
Gấp 14 lên 4 lần ta được kết quả là
\choice
{\True $56$}
{$46$}
{$66$}
{$54$}
\loigiai{
$14 \times 4 = 56$.\\
\textbf{Đáp án của câu hỏi là A}
}
\end{ex}

\begin{ex} %[3D4H24]
Gấp 18 lên 3 lần ta được kết quả là
\choice
{$48$}
{$58$}
{$52$}
{\True $54$}
\loigiai{
$18 \times 3 = 54$.\\
\textbf{Đáp án của câu hỏi là D}
}
\end{ex}

\begin{ex} %[3D4H24]
Số 12 gấp lên 6 lần thì bằng
\choice
{$62$}
{\True $72$}
{$82$}
{$68$}
\loigiai{
$12 \times 6 = 72$.\\
\textbf{Đáp án của câu hỏi là B}
}
\end{ex}

\begin{ex} %[3D4H24]
Đoạn thẳng $AB$ dài $8\text{ cm}$. Đoạn thẳng $CD$ dài gấp 4 lần đoạn thẳng $AB$. Đoạn thẳng $CD$ dài là
\choice
{$24\text{ cm}$}
{$36\text{ cm}$}
{\True $32\text{ cm}$}
{$28\text{ cm}$}
\loigiai{
$8 \times 4 = 32\text{ cm}$.\\
\textbf{Đáp án của câu hỏi là C}
}
\end{ex}

\begin{ex} %[3D4V24]
Bạn Nam có 9 viên bi, bạn Việt có số bi gấp 4 lần số bi của Nam. Hỏi cả hai bạn có tất cả bao nhiêu viên bi?
\choice
{\True $45$ viên}
{$36$ viên}
{$40$ viên}
{$48$ viên}
\loigiai{
Việt có: $9 \times 4 = 36$ viên bi.\\
Cả hai bạn có: $9 + 36 = 45$ viên bi.\\
\textbf{Đáp án của câu hỏi là A}
}
\end{ex}

\begin{ex} %[3D4V24]
Năm nay con 8 tuổi, tuổi bố gấp 5 lần tuổi con. Bố hơn con số tuổi là
\choice
{$30$ tuổi}
{$34$ tuổi}
{\True $32$ tuổi}
{$40$ tuổi}
\loigiai{
Tuổi bố: $8 \times 5 = 40$ tuổi.\\
Bố hơn con: $40 - 8 = 32$ tuổi.\\
\textbf{Đáp án của câu hỏi là C}
}
\end{ex}

\begin{ex} %[3D4V24]
Thùng thứ nhất đựng $15\text{ l}$ dầu, thùng thứ hai đựng gấp 3 lần thùng thứ nhất. Thùng thứ hai đựng số lít dầu là
\choice
{\True $45\text{ l}$}
{$35\text{ l}$}
{$50\text{ l}$}
{$40\text{ l}$}
\loigiai{
$15 \times 3 = 45\text{ l}$.\\
\textbf{Đáp án của câu hỏi là A}
}
\end{ex}

\begin{ex} %[3D4V24]
Một bao gạo nặng $12\text{ kg}$, một bao ngô nặng gấp 4 lần bao gạo. Hỏi cả hai bao nặng bao nhiêu ki-lô-gam?
\choice
{$48\text{ kg}$}
{$56\text{ kg}$}
{\True $60\text{ kg}$}
{$52\text{ kg}$}
\loigiai{
Bao ngô: $12 \times 4 = 48\text{ kg}$.\\
Cả hai bao: $12 + 48 = 60\text{ kg}$.\\
\textbf{Đáp án của câu hỏi là C}
}
\end{ex}

\begin{ex} %[3D4V24]
Số lớn gấp 5 lần số bé, biết số bé là 16. Tổng hai số là
\choice
{$80$}
{\True $96$}
{$90$}
{$86$}
\loigiai{
Số lớn: $16 \times 5 = 80$.\\
Tổng: $16 + 80 = 96$.\\
\textbf{Đáp án của câu hỏi là B}
}
\end{ex}

\begin{ex} %[3D4C24]
Nhà An nuôi 7 con gà, số vịt gấp 3 lần số gà. Sau đó mẹ mua thêm 5 con vịt nữa. Hỏi lúc này nhà An có tất cả bao nhiêu con vịt?
\choice
{$24$ con}
{$28$ con}
{$21$ con}
{\True $26$ con}
\loigiai{
Số vịt ban đầu: $7 \times 3 = 21$ con.\\
Số vịt lúc sau: $21 + 5 = 26$ con.\\
\textbf{Đáp án của câu hỏi là D}
}
\end{ex}

\textbf{Phần 2. Trắc nghiệm đúng sai}
\setcounterpageref{ex}{0}

\begin{ex} %[3D4H24]
Xét tính đúng sai về gấp một số lên nhiều lần:
\choiceTF
{\True Gấp 8 lên 3 lần được 24.}
{\True Gấp 15 lên 4 lần được 60.}
{Gấp 10 lên 5 lần được 15.}
{\True Gấp 10 lên 5 lần được 50.}
\loigiai{
\begin{itemize}
    \item $8 \times 3 = 24$: Đúng.
    \item $15 \times 4 = 60$: Đúng.
    \item Gấp là phép nhân ($10 \times 5 = 50 \neq 15$). Ý 3 Sai, Ý 4 Đúng.
\end{itemize}
\begin{itemize}
    \item a) \textbf{Đúng}.
    \item b) \textbf{Đúng}.
    \item c) \textbf{Sai}.
    \item d) \textbf{Đúng}.
\end{itemize}
}
\end{ex}

\begin{ex} %[3D4V24]
Tổ 1 trồng được 12 cây, tổ 2 trồng được số cây gấp 3 lần tổ 1. Xét tính đúng sai:
\choiceTF
{\True Tổ 2 trồng được 36 cây.}
{\True Cả hai tổ trồng được tất cả 48 cây.}
{Tổ 2 trồng nhiều hơn tổ 1 là 20 cây.}
{\True Tổ 2 trồng nhiều hơn tổ 1 là 24 cây.}
\loigiai{
\begin{itemize}
    \item Tổ 2: $12 \times 3 = 36$ cây. Đúng.
    \item Tổng: $12 + 36 = 48$ cây. Đúng.
    \item Hiệu: $36 - 12 = 24$ cây $\neq 20$. Ý 3 Sai, Ý 4 Đúng.
\end{itemize}
\begin{itemize}
    \item a) \textbf{Đúng}.
    \item b) \textbf{Đúng}.
    \item c) \textbf{Sai}.
    \item d) \textbf{Đúng}.
\end{itemize}
}
\end{ex}

\textbf{Phần 3. Trả lời ngắn}
\setcounterpageref{ex}{0}

\begin{ex} %[3D4H24]
Gấp 9 lên 6 lần ta được kết quả là bao nhiêu?
\shortans{54}
\loigiai{
\textbf{Đáp số:} $54$\\
$9 \times 6 = 54$.
}
\end{ex}

\begin{ex} %[3D4H24]
Gấp 24 lên 3 lần ta được kết quả là bao nhiêu?
\shortans{72}
\loigiai{
\textbf{Đáp số:} $72$\\
$24 \times 3 = 72$.
}
\end{ex}

\begin{ex} %[3D4V24]
Bạn Hoa có 8 quyển truyện, bạn Mai có số truyện gấp 4 lần bạn Hoa. Hỏi bạn Mai có bao nhiêu quyển truyện?
\shortans{32}
\loigiai{
\textbf{Đáp số:} $32$\\
$8 \times 4 = 32$ quyển truyện.
}
\end{ex}

\begin{ex} %[3D4V24]
Đoạn thẳng $AB$ dài $14\text{ cm}$, đoạn thẳng $CD$ dài gấp 5 lần đoạn thẳng $AB$. Tính độ dài đoạn thẳng $CD$ theo xăng-ti-mét.
\shortans{70}
\loigiai{
\textbf{Đáp số:} $70$\\
$14 \times 5 = 70\text{ cm}$.
}
\end{ex}

\begin{ex} %[3D4V24]
Một cửa hàng buổi sáng bán được $15\text{ kg}$ đường, buổi chiều bán gấp 4 lần buổi sáng. Hỏi cả hai buổi cửa hàng bán được bao nhiêu ki-lô-gam đường?
\shortans{75}
\loigiai{
\textbf{Đáp số:} $75$\\
Chiều: $60\text{ kg}$. Tổng: $15 + 60 = 75\text{ kg}$.
}
\end{ex}

\begin{ex} %[3D4C24]
Tìm một số biết rằng nếu gấp số đó lên 4 lần rồi trừ đi 15 thì được 45.
\shortans{15}
\loigiai{
\textbf{Đáp số:} $15$\\
Số đó $\times 4 = 45 + 15 = 60 \Rightarrow$ Số đó là $60 : 4 = 15$.
}
\end{ex}
